\documentclass[10pt,a4paper]{article}
\usepackage[T1]{fontenc}
\usepackage[utf8]{inputenc}
\usepackage[english]{babel}
\usepackage{amsmath,amssymb}
\usepackage{booktabs}
\usepackage[margin=2.0cm]{geometry}
\usepackage{microtype}

\setlength{\parindent}{0pt}
\setlength{\parskip}{0.45em}
\newcommand{\modelname}{\texttt{gerlee-history}}

\begin{document}

\begin{center}
  {\Large\bfseries Technical description of an adaptive historical-season regression model}\par
  \vspace{0.3em}
  {\modelname{} for the Swedish Influenza Forecast Hub}\par
  \vspace{0.2em}
  Philip Gerlee, Chalmers University of Technology
\end{center}

\section*{Purpose and data source}

The model produces point forecasts of weekly reported influenza cases for
Sweden (SE), Region Skåne (SE-M), and Region Västra Götaland (SE-O). Forecasts
are submitted for horizons $h=0,1,2,3$, where $h=0$ denotes the week ending on
the forecast reference date. The model uses only the time series published by
the forecast hub; no external epidemiological or demographic covariates are
included.

The central assumption is that influenza exhibits recurring seasonal shapes.
The number of cases in a given ISO week of the current season can therefore be
approximated by a linear combination of the counts in the corresponding week
of previous seasons. This combination is updated sequentially as new
observations become available during the current season.

An influenza season is defined as ISO week 30 in year $a$ through week 23 in
year $a+1$. Only historical seasons with complete observations for all three
locations are used. The 2020/21 season is excluded because the pandemic and
associated control measures produced highly atypical influenza circulation.
For the retrospective forecasts of 2025/26, eight seasons are used:
2016/17--2019/20 and 2021/22--2024/25.

\section*{Model formulation}

Consider one location at a time. Let $y_t$ denote the observed count in week
$t$ of the current season, and let
\[
  \mathbf{x}_t=(x_{t1},\ldots,x_{tn})^{\mathsf T}
\]
contain the counts in the corresponding ISO week of the $n$ historical
seasons. The model expectation is
\begin{equation}
  \widehat y_t = \mathbf{x}_t^{\mathsf T}\boldsymbol\beta,
  \qquad \boldsymbol\beta=(\beta_1,\ldots,\beta_n)^{\mathsf T},
  \qquad \beta_j\geq 0.
  \label{eq:forecast}
\end{equation}
At the beginning of the season, all historical seasons receive equal weight,
\begin{equation}
  \boldsymbol\beta_0=\frac{1}{n}\mathbf{1}.
  \label{eq:prior}
\end{equation}
The model starts updating the weights with observations from week 30. For a
forecast reference date $r$, let $\mathcal T_r$ contain all observed weeks from
week 30 through the Sunday preceding $r$. Let the rows of $\mathbf X_r$ be
$\mathbf{x}_t^{\mathsf T}$ and let $\mathbf y_r$ contain the corresponding
observations. The coefficients are estimated as
\begin{equation}
  \widehat{\boldsymbol\beta}_r=
  \underset{\boldsymbol\beta\geq0}{\operatorname{argmin}}
  \left\{
  \frac{\lVert\mathbf y_r-\mathbf X_r\boldsymbol\beta\rVert_2^2}{c^2}
  +\lambda\lVert\boldsymbol\beta-\boldsymbol\beta_0\rVert_2^2
  \right\}.
  \label{eq:ridge}
\end{equation}
Here, $c$ is the root mean square of the historical weekly values for the
location, calculated over weeks 30--23. This scaling gives $\lambda$ a
comparable interpretation across locations with different count levels. The
non-negativity constraint prevents hard-to-interpret compensation between
positive and negative seasonal weights. The coefficients are not constrained
to sum to one, allowing the model to adjust the overall level upward or
downward.

For reference date $r$, the four forecasts are
\begin{equation}
  \widehat y_{r+h}=\mathbf{x}_{r+h}^{\mathsf T}
  \widehat{\boldsymbol\beta}_r,\qquad h=0,1,2,3.
  \label{eq:horizons}
\end{equation}
The same coefficient vector is used for all four horizons, but it is
re-estimated for every forecast round. If the target ISO week is week 53,
which is absent from most historical seasons, the mean of weeks 52 and 1 is
used as the historical covariate.

\section*{Regularisation and calibration}

Equation~\eqref{eq:ridge} can be interpreted as a maximum-a-posteriori
estimate in a Gaussian regression model with a normal prior centred on the
equal weights in equation~\eqref{eq:prior}, truncated to the non-negative
orthant. Large values of $\lambda$ keep the model close to the historical
seasonal mean, whereas small values permit faster adaptation to the current
season.

One common regularisation value is used for all locations. The candidates
were
\[
  0.01,\ 0.03,\ 0.1,\ 0.3,\ 1,\ 3,\ 10,\ 30,\ 100,\ 300,\ 1000.
\]
They were compared using rolling forecast origins in earlier complete
seasons. For each calibration season, only still earlier seasons were used as
covariates, and at least three covariate seasons were required. The selection
criterion was mean absolute error across locations, forecast rounds, and the
four horizons. The best common value was
\[
  \boxed{\lambda=1}.
\]
The 2025/26 season was held out completely from this selection and was used
only for final retrospective evaluation.

\section*{Retrospective check and operational use}

For 2025/26, common regularisation gave an overall MAE of 96.7 cases,
compared with 98.5 when $\lambda$ was selected separately for each location.
MAE was 229.8 for Sweden, 27.6 for Skåne, and 32.8 for Västra Götaland. This
supports the simpler common specification, but it is evidence from a single
influenza season and should not be interpreted as a general performance
guarantee.

In operation, the model reads the data version available for the current
forecast round, removes observations after the Sunday preceding the reference
date, identifies complete previous seasons, and estimates
$\widehat{\boldsymbol\beta}_r$ separately for each location. Its output consists
of twelve non-negative point forecasts: three locations times four horizons.
The model does not quantify forecast uncertainty; probabilistic forecasts are
produced by the hub's separate ensemble model.

The main limitation is the assumption that future epidemic trajectories can
be represented by historical seasons. New viral variants, changes in testing,
or structural changes in reporting may therefore lead to systematic errors.
The fitted seasonal weights, their sum, and forecast errors should be
monitored during the season, but $\lambda$ must not be changed using outcomes
from a season that is being evaluated concurrently.

\end{document}
